\documentclass[11pt]{article}

\usepackage[utf8]{inputenc}
\usepackage[T1]{fontenc}
\usepackage{lmodern}
\usepackage{geometry}
\usepackage{amsmath,amssymb,amsfonts,amsthm,mathtools}
\usepackage{bm}
\usepackage{enumitem}
\usepackage{microtype}
\usepackage{hyperref}
\usepackage{titlesec}
\usepackage{booktabs}
\usepackage{array}
\usepackage{setspace}
\usepackage{mathrsfs}
\usepackage{physics}

\geometry{margin=1in}
\hypersetup{colorlinks=true, linkcolor=black, urlcolor=blue, citecolor=black}
\setlist[enumerate]{topsep=0.4em,itemsep=0.35em}
\setlist[itemize]{topsep=0.3em,itemsep=0.25em}
\titleformat{\section}{\large\bfseries}{\thesection.}{0.5em}{}
\titleformat{\subsection}{\normalsize\bfseries}{\thesubsection.}{0.5em}{}
\setstretch{1.06}

\newcommand{\Aether}{\AE{}ther}
\newcommand{\Aflow}{\AE{}ther-flow}
\newcommand{\TheoryName}{\Aflow{} Interpretation of Relativity}
\newcommand{\TheoryProgram}{\Aether{} / \Aflow{} framework}
\newcommand{\Sub}{\mathcal{A}}
\newcommand{\BridgeMetric}{\mathfrak{g}}
\newcommand{\Ell}{\mathcal{L}}

\newtheorem{definition}{Definition}
\newtheorem{proposition}{Proposition}
\newtheorem{remark}{Remark}
\newtheorem{corollary}{Corollary}
\newtheorem{theorem}{Theorem}

\title{\textbf{The \TheoryName{}}\\[0.4em]
\large Higher-Derivative Quartic Branch Same-Branch Renormalized Asymptotic Test}
\author{}
\date{}

\begin{document}

\maketitle

\begin{abstract}
The strict-branch decay/obstruction note isolated the first obstacle to a straightforward nonlinear-stability upgrade for the explicit minimal quartic-compatible completion \(S_{\mathrm{min},4}\): the maximal-gauge lapse perturbation \(n=N-1\) carries a generic Coulomb tail, and the scalar equation re-inserts that tail through the background-slope term \(\sigma_0 n\). That note deliberately stopped at the obstruction itself. The next disciplined question is narrower: before changing gauge, completion, or auxiliary architecture, can the same strict branch be repaired by renormalized asymptotic bookkeeping or a renormalized scalar variable?

This manuscript records that next step in a manuscript-safe form. It defines an exterior Coulomb profile
\[
n_{\mathrm{asy}}(t,x)=\chi_{\mathrm{ext}}(|x|)\frac{\mathfrak m_N(t)}{4\pi|x|},
\]
where \(\mathfrak m_N(t)\) is the lapse monopole already identified in the obstruction note, and it introduces a same-branch renormalized scalar
\[
\widetilde{\sigma}=\sigma-\sigma_0\Gamma_N,
\qquad
(\partial_t-\Ell_\beta)\Gamma_N=n_{\mathrm{asy}}.
\]
At the exact equation level this removes the explicit Coulomb forcing from the scalar transport law and replaces decay back to the unrenormalized flat target by decay relative to the modified asymptotic target \(\sigma_{\mathrm{asy}}=\sigma_0\Gamma_N\).

The note does \emph{not} prove that this repair already works. Instead it sharpens the same-branch burden to two concrete items: prove that the lapse admits a renormalized asymptotic decomposition \(n=n_{\mathrm{asy}}+n_{\mathrm{rem}}\) with integrable remainder on the wave zone, and prove that the corresponding correction profile can be commuted through the strict-branch corrected energy without destroying coercivity. The follow-on failure note \texttt{aether\_flow\_substrate\_higher\_derivative\_strict\_branch\_renormalized\_asymptotic\_failure.tex} now records that the specific first-pass scalar-variable package fixed here does not satisfy that second requirement.
\end{abstract}

\section{Introduction}

The higher-derivative quartic branch now has a fixed explicit completion, a reduced Einstein--quartic-scalar \(3+1\) system, a first local well-posedness theorem, a corrected coercive \(T\sim \varepsilon^{-1}\) bootstrap, a same-branch strict-elliptic relaxation test, and a first obstruction note for the nonlinear-stability upgrade problem. The active issue is therefore not whether a branch exists or whether the first obstruction is still hidden. It has already been identified precisely: on the present strict auxiliary elliptic-control branch, the lapse perturbation develops a Coulombic tail and feeds back into the scalar equation through \(\sigma_0(N-1)\).

That identification changes the next step again. The immediate burden is no longer to ask vaguely for ``more decay,'' and it is not yet to abandon the current branch formulation. The first honest next test is whether the obstruction can be removed \emph{within the same strict branch} by explicitly tracking and subtracting the long-range asymptotic profile generated by the lapse monopole.

\paragraph{Framework and claim boundary.}
\TheoryProgram{} denotes the broader \Aether{} / \Aflow{} framework built on the claims that the \Aether{} is the underlying four-dimensional substrate of reality, the \Aflow{} is its intrinsic ordered motion, observed three-dimensional space is the local experiential slice of that deeper substrate, S-time is the experienced order of change arising from matter, light, and the \Aflow{}, and observed expansion is the three-dimensional appearance of deeper four-dimensional ordered motion. Within that framework, \TheoryName{} denotes the exact-closure theory in which the effective gravitational dynamics are adopted to be exactly Einsteinian with universal matter coupling. In the active sequence, \emph{adoption} means use of that established relativistic dynamics without claiming substrate derivation, while \emph{derivation} is reserved for a first-principles recovery from explicit substrate variables.

The present note stays within that boundary. It does not claim that the long-range lapse obstruction has already been removed. It records the narrowest same-branch renormalized asymptotic test that must be attempted before the active sequence turns to a different gauge, completion, or auxiliary architecture.

\section{The Same-Branch Renormalization Question}

The obstruction note showed that the scalar transport law on the present strict branch has the form
\begin{equation}
(\partial_t-\Ell_\beta)\sigma
=
N\frac{\pi_\sigma}{m_S^2}
+ \sigma_0 n,
\qquad
n\equiv N-1,
\label{eq:sigmaeq}
\end{equation}
while the lapse perturbation solves an elliptic equation whose generic far field is
\begin{equation}
n(t,x)=\frac{\mathfrak m_N(t)}{4\pi|x|}+\mathcal O(|x|^{-2}).
\label{eq:coulombtail}
\end{equation}
The explicit obstruction came from integrating the \(t^{-1}\) wave-zone forcing produced by the leading term in \eqref{eq:coulombtail}.

\begin{definition}[Exterior Coulomb profile]
Fix a smooth radial cutoff \(\chi_{\mathrm{ext}}\in C^\infty(\mathbb R^3)\) satisfying
\[
0\le \chi_{\mathrm{ext}}\le 1,
\qquad
\chi_{\mathrm{ext}}(x)=0 \text{ for } |x|\le 1,
\qquad
\chi_{\mathrm{ext}}(x)=1 \text{ for } |x|\ge 2.
\]
The \emph{same-branch exterior Coulomb profile} is
\begin{equation}
n_{\mathrm{asy}}(t,x)
\equiv
\chi_{\mathrm{ext}}(x)\frac{\mathfrak m_N(t)}{4\pi|x|},
\label{eq:nasy}
\end{equation}
where \(\mathfrak m_N(t)\) is the lapse monopole from the strict-branch obstruction note.
\end{definition}

\begin{remark}
Definition 1 does not modify the completion, the gauge, or the auxiliary equations. It only isolates the explicit long-range term already forced by the same strict-branch lapse equation.
\end{remark}

\begin{definition}[Renormalized scalar corrector and renormalized scalar]
Let \(\Gamma_N\) be any scalar solving
\begin{equation}
(\partial_t-\Ell_\beta)\Gamma_N
=
n_{\mathrm{asy}},
\qquad
\Gamma_N|_{t=0}=0.
\label{eq:Gamma}
\end{equation}
Define the \emph{renormalized same-branch asymptotic target}
\begin{equation}
\sigma_{\mathrm{asy}}\equiv \sigma_0\Gamma_N
\label{eq:sigmaasy}
\end{equation}
and the \emph{renormalized scalar variable}
\begin{equation}
\widetilde{\sigma}
\equiv
\sigma-\sigma_{\mathrm{asy}}
=
\sigma-\sigma_0\Gamma_N.
\label{eq:sigmatilde}
\end{equation}
\end{definition}

\begin{proposition}[Exact renormalized scalar identity]
The renormalized scalar variable \(\widetilde{\sigma}\) satisfies
\begin{equation}
(\partial_t-\Ell_\beta)\widetilde{\sigma}
=
N\frac{\pi_\sigma}{m_S^2}
+ \sigma_0\bigl(n-n_{\mathrm{asy}}\bigr).
\label{eq:rensigma}
\end{equation}
\end{proposition}

\begin{proof}
Subtract \(\sigma_0\) times \eqref{eq:Gamma} from \eqref{eq:sigmaeq} and use Definition 2.
\end{proof}

\begin{remark}
Equation \eqref{eq:rensigma} is the precise same-branch reason to try a renormalized scalar variable before changing formulation. It removes the explicit Coulomb profile algebraically, with no new gauge condition and no new nonlinear completion.
\end{remark}

\section{Renormalized Asymptotic Package}

The renormalized variable is useful only if the remainder \(n-n_{\mathrm{asy}}\) decays better than the raw lapse perturbation. That is the next mathematical burden.

\begin{definition}[Renormalized asymptotic package]
Fix a strict-branch solution satisfying the linear-compatible decay package of the strict-branch obstruction note. The solution is said to satisfy the \emph{renormalized asymptotic package} if there exist constants \(C_{\mathrm{ren}}>0\) and \(\delta_{\mathrm{ren}}>0\) such that:
\begin{enumerate}
    \item the lapse admits the decomposition
    \begin{equation}
    n=n_{\mathrm{asy}}+n_{\mathrm{rem}}
    \label{eq:ndecomp}
    \end{equation}
    with \(n_{\mathrm{asy}}\) from Definition 1
    \item the wave-zone remainder is integrable in the sense that
    \begin{equation}
    \sup_{|x|\simeq 1+t}|n_{\mathrm{rem}}(t,x)|
    \le
    C_{\mathrm{ren}}\varepsilon^2(1+t)^{-1-\delta_{\mathrm{ren}}}
    \label{eq:nrem}
    \end{equation}
    for all \(t\ge 0\)
    \item the monopole is asymptotically controlled by
    \begin{equation}
    |\mathfrak m_N(t)|
    +(1+t)^{1+\delta_{\mathrm{ren}}}|\partial_t\mathfrak m_N(t)|
    \le
    C_{\mathrm{ren}}\varepsilon^2.
    \label{eq:mcontrol}
    \end{equation}
\end{enumerate}
\end{definition}

\begin{remark}
Definition 3 is intentionally a same-branch test package, not a theorem. Item \eqref{eq:nrem} is the exact strengthening needed to replace the logarithmically divergent forcing in the obstruction note by an \(L_t^1\) term on the wave zone. Item \eqref{eq:mcontrol} is the minimum asymptotic bookkeeping needed to track the renormalized profile itself without introducing an uncontrolled time-dependent coefficient.
\end{remark}

\begin{proposition}[Wave-zone forcing becomes integrable after renormalization]
Assume the renormalized asymptotic package. Then the forcing term in \eqref{eq:rensigma} obeys
\begin{equation}
\int_0^\infty
\sup_{|x|\simeq 1+t}
\left|
\sigma_0\bigl(n-n_{\mathrm{asy}}\bigr)(t,x)
\right|dt
\le
C_{\mathrm{ren}}|\sigma_0|\varepsilon^2
\int_0^\infty (1+t)^{-1-\delta_{\mathrm{ren}}}\,dt
<\infty.
\label{eq:integrableforcing}
\end{equation}
Hence the specific logarithmic divergence identified in the strict-branch obstruction note is absent from the renormalized scalar equation.
\end{proposition}

\begin{proof}
Insert the decomposition \eqref{eq:ndecomp} into \eqref{eq:rensigma}. Estimate the wave-zone remainder using \eqref{eq:nrem}. The integral in \eqref{eq:integrableforcing} converges because \(\delta_{\mathrm{ren}}>0\).
\end{proof}

\begin{remark}
If \(\mathfrak m_N(t)\) does not decay integrably, then \(\Gamma_N\) carries the secular exterior profile generated by the accumulated monopole \(\mathfrak M_N(t)=\int_0^t\mathfrak m_N(s)\,ds\). On the wave zone this appears as a modified asymptotic target rather than decay back to the unrenormalized flat state. The role of \(\widetilde{\sigma}\) is precisely to absorb that long-range phase into the asymptotic target rather than to pretend that the unrenormalized variable should decay directly to zero.
\end{remark}

\section{Decision Criterion on the Present Branch}

The preceding section isolates the same-branch repair in a precise way. What remains is to decide whether the strict branch can actually support that repair quantitatively.

\begin{theorem}[What must be shown before leaving the strict branch]
The first obstruction identified in the strict-branch decay/obstruction note is removed \emph{within the same strict branch} if the following two additional items can be proved for the explicit minimal quartic-compatible completion \(S_{\mathrm{min},4}\):
\begin{enumerate}
    \item the renormalized asymptotic package of Definition 3
    \item a renormalized high-order energy estimate for \(\widetilde{\sigma}\) and the remaining strict-branch variables that preserves the corrected coercive structure already recorded for the unrenormalized system.
\end{enumerate}
If either item fails, only then is it mathematically disciplined for the active sequence to turn next to a different gauge, completion, or auxiliary architecture.
\label{thm:decision}
\end{theorem}

\begin{proof}
Proposition 2 shows that Definition 3 removes the specific nonintegrable wave-zone forcing isolated in the earlier obstruction note. Thus the first obstruction is no longer the raw Coulomb tail itself, but whether that tail can be isolated sharply enough and carried through the commuted coercive energy without introducing a new nonintegrable loss. If both items in the theorem hold, then the previously identified obstacle has been repaired on the same branch rather than by changing formulation. If one of them fails, the same-branch renormalized route has failed at its narrowest honest formulation, and a new gauge, completion, or auxiliary architecture becomes the next disciplined option.
\end{proof}

\begin{corollary}[Immediate manuscript priority]
At current scope, the next manuscript burden is not yet a different gauge, not yet a different completion, and not yet a different auxiliary architecture. It is the proof-or-failure test singled out in Theorem \ref{thm:decision}: establish or rule out the renormalized asymptotic package and its compatible renormalized energy bookkeeping on the present strict branch.
\end{corollary}

\begin{proof}
This is a direct restatement of Theorem \ref{thm:decision}.
\end{proof}

\section{What This Step Establishes}

The present manuscript establishes the following.
\begin{enumerate}
    \item It records the narrowest same-branch renormalized response to the long-range lapse obstruction already isolated in the previous note.
    \item It introduces an explicit renormalized scalar variable \(\widetilde{\sigma}\) and a matching asymptotic target \(\sigma_{\mathrm{asy}}\) without changing the strict branch, the maximal-gauge formulation, or the explicit completion.
    \item It identifies the exact asymptotic input needed for that repair to work: a Coulomb-profile decomposition of the lapse with integrable wave-zone remainder.
    \item It sharpens the active decision tree: prove or refute that same-branch renormalized package first, and only after failure turn to a different gauge, completion, or auxiliary architecture.
\end{enumerate}

It does not establish the following.
\begin{enumerate}
    \item It does not prove that the renormalized asymptotic package already holds for the current strict branch.
    \item It does not yet prove a renormalized global nonlinear-stability theorem.
    \item It does not change the explicit completion \(S_{\mathrm{min},4}\), the maximal-gauge reduction, or the strict auxiliary elliptic-control hypotheses.
    \item It does not rule out future need for a different gauge, completion, or auxiliary architecture if the same-branch renormalized test fails.
\end{enumerate}

\section{Conclusion}

The strict-branch obstruction note identified the first problem honestly: the lapse monopole generates a Coulombic tail, and the unrenormalized scalar variable sees that tail as a nonintegrable forcing through \(\sigma_0(N-1)\). The right next move is now equally narrow and honest. Before redesigning the formulation, the active sequence should test whether the same strict branch can absorb that tail into a renormalized asymptotic target and a renormalized scalar variable.

This note makes that test precise. The immediate burden is now to prove or refute the renormalized asymptotic package and its compatible high-order bookkeeping on the present branch. Only if that same-branch route fails does the active program point next to a different gauge, completion, or auxiliary architecture.

\end{document}
